\documentclass[10pt,a4paper]{amsart}
\usepackage[margin=19mm]{geometry}
\usepackage{amsmath,amssymb,mathtools}
\usepackage[scr=rsfs,cal=euler]{mathalfa}
\usepackage{xfrac}
\usepackage[unicode,hidelinks,bookmarks=false]{hyperref}
\newcommand{\ie}{i.e.\ }
\newcommand{\cf}{cf.\ }
\newcommand{\resp}{respectively\ }
\newcommand{\Subsection}{\subsection}
\providecommand{\nobreakdash}{\nobreak\hbox{-}}
\setlength{\emergencystretch}{2em}
\multlinegap0pt
\usepackage{amssymb}
\usepackage{xcolor}
\usepackage[normalem]{ulem}
\usepackage{tikz}
\newtheorem{lemma}{Lemma}
\newtheorem{remark}{Remark}
\title{Rings whose Non-Units are a Unit Multiple \\ of an Element from $\sqrt{\Delta(R)}$}
\author{Omid Hasanzadeh, Ahmad Moussavi, Peter Danchev}
\date{Source: arXiv:2602.14600v1 (CC BY 4.0)}
\begin{document}
\maketitle

\noindent\emph{Source and licence.} Rings whose Non-Units are a Unit Multiple \\ of an Element from $\sqrt{\Delta(R)}$. Version arXiv:2602.14600v1 (CC BY 4.0), distributed by arXiv under the Creative Commons Attribution 4.0 International licence, http://creativecommons.org/licenses/by/4.0/, as recorded in the arXiv licence metadata for this exact version and shown on the abstract page https://arxiv.org/abs/2602.14600v1. Full licence notice: \texttt{licenses/arxiv-2602.14600-CC-BY-4.0.txt}.\par\smallskip

\noindent\textit{Editorial scope.} Counted unit: the article's lemma lemma (source0.tex lines 676--682). The article's complete original proof of it is delivered with it (source0.tex lines 684--698, one \texttt{proof} environment of 15 lines). No mathematics is written, repaired or re-derived: the statement and the proof are the article's own lines, in the article's own notation, and no retained line is substituted or re-flowed.  Retained uncounted as the source-local context the proof needs: lines 247--271; lines 305--330.  Nothing was omitted from any retained span, so the package carries no omission record.  No retained line contains a citation, so the wrapper carries no bibliography and no bibliography entry.  No retained line contains a figure environment, an image-inclusion command or drawing code, so no image is delivered and the package declares no asset dependency.  Editorial formatting only, in addition: an AMS \texttt{amsart} preamble that restates the article's own declarations and macros used by the retained lines, namely the environments \texttt{lemma}, \texttt{remark} on separate counters, together with the standard packages those lines need.  The retained environments print in this document as Lemma 1, Remark 1 in order of appearance, whereas the article prints the same items with the numbering of its own full document.  The complete native source of the article as distributed in the arXiv e-print for this exact version is bundled unchanged as \texttt{sources/194726/source0.tex}.
\begin{lemma}\label{property of J sharp}
For each ring \( R \), we have:

(1) If \( a \in \sqrt{\Delta(R)} \) and \( b \in U(R) \) with \( ab = ba \), then \( ab \in \sqrt{\Delta(R)} \).

(2) If \( a \in \sqrt{\Delta(R)} \), then \( 1 - a \in U(R) \).

(3) If \( a \in \sqrt{\Delta(R)} \cap C(R) \), then \( a \in \Delta(R) \).

(4) For every ideal \( I \subseteq J(R) \), we have \( \sqrt{\Delta(R/I)} = \sqrt{\Delta(R)}/I \).

(5) If \( R = \prod_{i=1}^k R_i \), then \( \sqrt{\Delta(R)} = \prod_{i=1}^k \sqrt{\Delta(R_i)} \).

(6) $\sqrt{\Delta(R)} \Delta(R) \subseteq \Delta(R)$.

(7) $\sqrt{\Delta(R)} + \Delta(R) \subseteq \sqrt{\Delta(R)}$.

(8) If $a\in U(\sqrt{\Delta}(R))\cap C(R)$, then $1-a\in U(R)$.

(9) If $a\in U(\sqrt{\Delta}(R))$, then, for each $u,v\in U(R)$, $uav\in U\sqrt{\Delta}(R)$.

(10) If $a\in \sqrt{\Delta(R)}$, then, for each $u\in U(R)$, $u^{-1}au\in \sqrt{\Delta(R)}$.

(11) If $ac\in\sqrt{\Delta(R)}$, where $c\in C(R)\cap U(R)$, then $a\in \sqrt{\Delta(R)}$.
\end{lemma}

\begin{remark}\label{local}
For each ring \( R \), we have validity of the following four points:

(1) It is well known that $\Delta(R)\subseteq \sqrt{\Delta(R)}$ and ${\rm Nil}(R)\subseteq \sqrt{\Delta(R)}$;
however, the converse implications do {\it not} hold in general. To check that, set $S:=M_{2}(\mathbb{F}_{2})$,
$K=\mathbb{F}_{2}[[x]]$, and $L=S\times K$. Consider the element
$
A=\begin{pmatrix}
	1 & 1\\
	1 & 1
\end{pmatrix}\in S.
$
Then, one verifies that $A^{2}=0\in \Delta(S)$, which yields that $A\in \sqrt{\Delta(S)}$, while
$A\not\in \Delta(S)$. Moreover, $x\in J(K)\subseteq \sqrt{\Delta(K)}$, but
$x\not\in {\rm Nil}(K)$. Now, let $a:=(A,x)\in L$. Utilizing Lemma\ref{property of J sharp}(5), we obtain
$a\in \sqrt{\Delta(L)}$; however, $a\not\in J(L)$ and $a\not\in {\rm Nil}(L)$.

(2) \( \sqrt{\Delta(R)} \cap {\rm Id}(R) = \{0\} \). In fact, if \( e \in \sqrt{\Delta(R)} \cap {\rm Id}(R) \), then Lemma~\ref{property of J sharp}(2) ensures that \( 1-e \in {\rm Id}(R) \cap U(R) \), and so \( e=0 \).

(3) \( R \) is a local ring if, and only if, \( R=U(R) \cup \sqrt{\Delta(R)} \). Indeed, if \( R \) is local, then
one sees that \[ R=U(R) \cup J(R) \subseteq U(R) \cup \sqrt{\Delta(R)} \subseteq R ,\] and thus \( R=U(R) \cup \sqrt{\Delta(R)} \).

Conversely, suppose \( R=U(R) \cup \sqrt{\Delta(R)} \) and \( a \notin U(R) \). Then, \( a \in \sqrt{\Delta(R)} \), so that Lemma~\ref{property of J sharp}(2) employs to get that \( 1-a \in U(R) \), whence \( R \) is a local ring.

(4) \( \sqrt{\Delta(R)} \cap U(R) = \emptyset \). Indeed, if \( u \in \sqrt{\Delta(R)} \cap U(R) \), then there exists \( n \in \mathbb{N} \) such that \( u^n \in \Delta(R) \), and therefore \( 0=1-u^{-n}u^n \in U(R) \), which is an obvious contradiction.
\end{remark}

\begin{lemma}
Let $R$ be a ring. Then,

(1) $R$ is uniquely $U\sqrt{\Delta}$ if and only if $R$ is a division ring.

(2) $R$ is strongly $U\sqrt{\Delta}$ if and only if $R$ is local.
\end{lemma}

\begin{proof}
(1) Suppose $d \in \sqrt{\Delta(R)}$. Then, we may write
\[
			d = 1 \cdot d = (1 - d)(1 - d)^{-1}d.
\]
Since $d$ and $(1 - d)^{-1}$ commute one another, Lemma \ref{property of J sharp}(1) is working to get $(1 - d)^{-1}d \in \sqrt{\Delta(R)}$, and hence
	\[
			1 = 1 - d \Rightarrow d = 0.
	\]
Therefore, $\sqrt{\Delta(R)} = 0$, and so $R$ is a division ring. The converse is rather clear, so we omit details.
			
(2) Suppose $R$ is strongly $U\sqrt{\Delta}$. Let $x \notin U(R)$. Then, $x = ud = du$ for some $u \in U(R)$ and $d \in \sqrt{\Delta(R)}$. Playing with Lemma \ref{property of J sharp}(1), we arrive at $x \in \sqrt{\Delta(R)}$, and thus $R$ is local in virtue of Remark \ref{local}(3).
			
Conversely, assume $R$ is local. Thus, again in view of Remark \ref{local}(3), all non-units are in $\sqrt{\Delta(R)}$. Letting $d \in \sqrt{\Delta(R)}$, then $d = 1 \cdot d$ is a strongly $U\sqrt{\Delta}$ decomposition. Finally, $R$ is strongly $U\sqrt{\Delta}$, as claimed.
\end{proof}

\end{document}
